\section{Background and Motivation}

\paragraph{Previous Solutions.}
Knowledge distillation (KD) typically trains a smaller student model to imitate the predictive behavior of a larger or more capable teacher~\citep{hinton2015distilling_kd,ko2024distillm_opd}. For autoregressive language models, common offline approaches include sequence-level distillation, in which the student is trained on teacher-generated sequences~\citep{kim2016sequence}, and supervised token-level KD, in which the student matches the teacher's next-token distributions along prefixes induced by a fixed offline set of output sequences~\citep{agarwal2024onpolicy_opsd}. These output sequences may be ground-truth or teacher-generated. In their purely offline form, both approaches optimize over an output-prefix distribution that does not track the current student's rollout distribution. During autoregressive inference, however, the student conditions on its own previous predictions. A deviation can therefore shift subsequent prefixes away from those observed during offline distillation, creating a training--inference distribution mismatch and allowing errors to compound across decoding steps~\citep{lin2020imitkd,agarwal2024onpolicy_opsd}.


\paragraph{On-Policy Distillation Mechanism.}
On-policy distillation (OPD) collects teacher supervision on outputs generated by the current student, rather than relying only on fixed offline sequences. During training, the student generates responses or trajectories, and the teacher provides supervision on the resulting prefixes or states. The student is then updated, and new rollouts are generated with the updated policy. In a common white-box setting, the teacher provides next-token distributions on student-generated prefixes, and the student minimizes a token-level divergence from the teacher~\citep{agarwal2024onpolicy_opsd}. The defining property of OPD is therefore that supervision is collected on states induced by the current student, rather than the use of any particular distillation objective~\citep{lu2025onpolicydistillation,song2026survey}.

Training on student-generated states helps reduce the mismatch between offline training data and the prefixes encountered during autoregressive inference, while providing guidance on errors made by the student itself~\citep{agarwal2024onpolicy_opsd,ko2024distillm_opd}. However, OPD also turns distillation into an iterative pipeline that must coordinate student rollout, teacher inference, supervision construction, and model optimization. Since OPD methods differ in teacher access, feedback form, and training objectives~\citep{song2026survey}, a practical framework should support this shared pipeline while allowing method-specific supervision to be implemented consistently.
